\subsection{The three subtasks}
\label{sec:task-subtasks}

The Touch\'{e} 2026 lab~\cite{heinrich2026touche} (task description and labels at \url{https://touche.webis.de/clef26/touche26-web/fallacy-detection.html}) defines three subtasks over the Reddit informal-fallacy data of Sahai et al.~\cite{sahai2021invisiblewall}, with evaluation hosted on TIRA~\cite{frobe2023tira}. The 2026 edition is the first to feature fallacy detection as a dedicated lab task, continuing a Touch\'{e} lineage that has progressively shifted from argument retrieval~\cite{bondarenko2022touche,bondarenko2023touche} to full argumentation systems~\cite{kiesel2024touche,kiesel2025touche}. ST1 is a binary decision on whether a given comment, presented together with its parent comment and the title of the originating thread, commits any informal fallacy. ST2 conditions on a positive ST1 decision and asks for the specific fallacy type from an eight-way inventory: \textsf{authority}, \textsf{black\_white}, \textsf{hasty\_generalization}, \textsf{natural}, \textsf{population}, \textsf{slippery\_slope}, \textsf{tradition}, \textsf{worse\_problems}. ST3 conditions on a negative ST1 decision and asks for the argument scheme of the comment from a four-way inventory operationalised by the lab task description~\cite{heinrich2026touche}, whose conceptual antecedents are the Walton--Reed--Macagno scheme catalogue~\cite{walton2008schemes} and the argumentation-profile framing of Macagno~\cite{macagno2022argumentation}: \textsf{pi} (practical-internal), \textsf{pe} (practical-external), \textsf{ei} (epistemic-internal), \textsf{ee} (epistemic-external). The \textsf{pi}/\textsf{pe}/\textsf{ei}/\textsf{ee} abbreviations are the labels used throughout our internal logs and the rest of this notebook. The label set positions naturally inside the broader fallacy taxonomies surveyed by Helwe et al.~\cite{helwe2024mafalda}, although we do not retrain or evaluate on MAFALDA. Evaluation is by macro-averaged F1 in all three subtasks, which puts roughly equal weight on each label and turns rare-class behaviour into a binding constraint on the headline number --- a point we return to in \S\ref{sec:disc-pe}.

\subsection{Base and enhanced tracks}
\label{sec:task-tracks}

Each subtask is evaluated in two tracks. The \emph{base} track exposes only original Reddit fields --- in our setup, \texttt{text\_raw}, \texttt{text\_raw\_parent}, \texttt{text\_raw\_title}, \texttt{text\_base}, \texttt{argument\_base} --- which the model concatenates into a single input. The \emph{enhanced} track additionally exposes two organiser-supplied paraphrases, \texttt{text\_enhanced} and \texttt{argument\_enhanced}, that were rewritten with knowledge of the gold labels in order to make the argumentative structure of the comment more explicit. We flag the rewriting-with-label-knowledge here because it directly affects how the enhanced-track numbers should be read; we defer the quantitative diagnostic and the defensive framing to \S\ref{sec:disc-leakage}, where the four-part argument can be made in one place. The lab permits up to three runs per (subtask, track), for a maximum of nine submissions; we used six (one per cell), populating one slot per subtask per track and not using the additional two-runs-per-track allowance.

\subsection{Dataset statistics}
\label{sec:task-stats}

% TODO senior: ST1/ST2 dev sizes
Following the official Touch\'{e} 2026 release, we hold out a stratified development partition from the published training data; the ST3 partition contains 473 examples. % from exp_0124
For ST1 and ST2 we adopt the analogous stratified split; exact sizes are deferred to the camera-ready once the test phase concludes. The ST3 partition is the one we use for the 5-fold cross-validation runs reported in \S\ref{sec:results}. % from exp_0124 exp_0153
Class distributions follow the underlying Sahai et al.\ data and are long-tailed; the ST3 \textsf{pe} class in particular is small enough to be the binding constraint on macro-F1, which we revisit in \S\ref{sec:thresholds} and \S\ref{sec:disc-pe}.
